\section{관련 연구}
{\AA}str\"om과 H\"agglund는 PID 튜닝, 잡음, 강인성, 구현을 다룬다 \cite{astrom2006}. 안티와인드업 추적과 조건부 적분은 \cite{astrom1989,rundqwist1990}에 정립되어 있으며, 이는 실제 액추에이터 한계를 제어기 상태에 반영해야 함을 뒷받침한다. Anderson과 Moore는 Riccati 및 선형 이차 제어의 기초를 제시하고 \cite{anderson1990}, Zhou, Doyle, Glover는 공칭 성능, 강인 성능, 강인 안정성을 구분한다 \cite{zhou1996}. 최근 열유동 비교 연구 역시 명시적 튜닝 목적함수의 중요성을 보인다 \cite{teppa2025}. 영차 유지의 정확한 구간 입력 해석과 전진 Euler 이산 LQI 규약은 권위 있는 기술 자료를 따른다 \cite{mathworks_zoh,mathworks_lqi}. 이 검토는 포괄적인 신규성 조사가 아니다.
